\section{HabitatNet Architecture}

\label{sec:architecture}
% ══════════════════════════════════════════════════════════════════════════════

\subsection{Spatial Graph Construction}

\paragraph{Grid resolution.} We discretize the study area into 1\,km
hexagonal cells, following \citet{jetz2019essential}. Hexagonal grids
minimize edge effects and provide uniform adjacency
\citep{birch2007rectangular}.

\paragraph{Node features.} Each cell's feature vector $\mathbf{x}_i \in
\R^{42}$ comprises:
\begin{itemize}
  \item 19 bioclimatic variables from WorldClim v2.1
    \citep{fick2017worldclim}
  \item 5 topographic variables: elevation, slope, aspect,
    topographic roughness index (TRI), topographic wetness index (TWI)
  \item 10 land cover class fractions from ESA WorldCover
    \citep{zanaga2022esa}
  \item 4 vegetation indices from MODIS: NDVI mean, NDVI seasonality,
    LAI, FPAR
  \item 4 hydrological variables: distance to water, flow accumulation,
    precipitation seasonality, aridity index
\end{itemize}

\paragraph{Edge construction.} Edges connect each cell to its six
hexagonal neighbors. Long-range edges (up to 50\,km) are added for cells
along waterways and ridgelines, reflecting known dispersal corridors.

\paragraph{Edge weights.} Dispersal permeability $w_{ij}$ is initialized
from land cover resistance values but is subsequently refined through
learning (Section~\ref{sec:edge_learning}).

\subsection{Network Architecture}

HabitatNet consists of four modules:

\begin{enumerate}
  \item \textbf{Feature Encoder}: A shared 3-layer MLP that projects raw
    environmental features into a latent space:
    \begin{equation}
      \mathbf{h}_i^{(0)} = \text{MLP}_{\text{enc}}(\mathbf{x}_i)
      \in \R^{128}
    \end{equation}

  \item \textbf{Graph Message Passing}: $L$ layers of graph attention
    \citep{velickovic2018graph} with edge-conditioned convolution:
    \begin{equation}
      \mathbf{h}_i^{(\ell+1)} = \sigma\!\left(
        \sum_{j \in \mathcal{N}(i)} \alpha_{ij}^{(\ell)} \,
        \mathbf{W}^{(\ell)} \mathbf{h}_j^{(\ell)} +
        \mathbf{W}_e^{(\ell)} \mathbf{e}_{ij}
      \right)
    \end{equation}
    where $\alpha_{ij}^{(\ell)}$ are learned attention coefficients,
    $\mathbf{e}_{ij}$ is the edge feature vector, and
    $\sigma$ is the GELU activation.

  \item \textbf{Species-Specific Readout}: For each species $k$, an
    attention-weighted readout layer computes habitat suitability:
    \begin{equation}
      p_i^{(k)} = \sigma_{\text{sig}}\!\left(
        \mathbf{w}_k^\top \mathbf{h}_i^{(L)} + b_k
      \right)
    \end{equation}
    where $\mathbf{w}_k \in \R^{128}$ and $b_k \in \R$ are
    species-specific parameters.

  \item \textbf{Connectivity Module}: A spectral graph convolution layer
    that computes Fiedler values \citep{fiedler1973algebraic} of the
    predicted habitat subgraph, enabling differentiable connectivity
    optimization.
\end{enumerate}

\subsection{Edge Weight Learning}
\label{sec:edge_learning}

Rather than relying on hand-crafted resistance surfaces, HabitatNet
learns edge weights through a parameterized function:
\begin{equation}
  w_{ij} = \sigma_{\text{sig}}\!\left(
    \text{MLP}_{\text{edge}}([\mathbf{x}_i \| \mathbf{x}_j \|
    \mathbf{g}_{ij}])
  \right)
\end{equation}
where $\mathbf{g}_{ij}$ encodes geographic distance and elevation
difference, and $\|$ denotes concatenation. This allows the model to
discover species-group-specific resistance patterns from occurrence data.

\subsection{Climate Projection Integration}

For future projections, we replace the 19 bioclimatic variables with
downscaled CMIP6 projections \citep{eyring2016overview}. We use the
ensemble median of 12 general circulation models (GCMs) under four SSP
scenarios. The remaining features (topography, land cover) are kept
static for the base analysis but can be updated with projected land-use
change scenarios from LUH2 \citep{hurtt2020harmonization}.

To account for GCM uncertainty, we propagate the full ensemble through
HabitatNet and report the median and interquartile range of suitability
predictions.


% ══════════════════════════════════════════════════════════════════════════════
